\section{A direct quartic from one certificate for two front signs}
\label{sec:direct-front}

The second contribution is a local arithmetic simplification of the direct
C222 construction.  Its principal equation has 248 unrestricted integer
variables, degree four, 5,321 collected monomials, and coefficient height
128,446.  It is the sum of the squares of the 182 explicit quadratic
residuals in
\texttt{equations/integer248/equation.json}.  Here and below, coefficient
height is measured after collecting the entire polynomial.

\begin{center}
\begin{tabular}{lrrrrr}
Profile & Variables & Residuals & Degree & Terms & Height\\ \hline
Common Remainders parent & 252 & 182 & 4 & 5,340 & 128,446\\
Principal & 248 & 182 & 4 & 5,321 & 128,446\\
Balanced projection & 246 & 180 & 4 & 5,048 & 195,955\\
Dimension projection & 245 & 179 & 4 & 21,845 & 314,701
\end{tabular}
\end{center}

These are three achieved projections of one new raw source, not an
assertion that every better projection has been excluded.  The principal
profile strictly improves the dimension and term count of its matched
parent while preserving degree, residual count, and coefficient height.
The new mathematical step is a proof that one product certificate can
replace two front sign certificates, even though the product alone has
an unwanted negative branch.  The exact factorial digit already present
in the source excludes that branch.

\subsection{Literal source change and inherited interface}

Put
\[
 G(a,b,c)=a^2+b^2+c(c+1).
\]
This is nonnegative on $\Z^3$ and represents every nonnegative integer.
Indeed, the three-square theorem represents $4n+1$ as three squares; its
residues modulo four put exactly one odd square in that representation.
Writing it as $(2c+1)^2$ and the others as $(2a)^2,(2b)^2$ gives $n=G(a,b,c)$.

Use the following semantic names for the unchanged source coordinates:
\[
\begin{split}
 r&=\texttt{fact0.remainder},\quad
 s=\texttt{fact0.slack},\quad F=\texttt{fact0.output},\\
 B&=\texttt{choose0.output},\quad
 V=\texttt{pow0.output},\quad q=\texttt{pow2.output}.
\end{split}
\]
The old front separately imposed $r\geq0$, $s\geq0$, and
$o\geq0$, where $q=B+o+1$.  Replace its two first ternary guards by
\begin{equation}
 (r+1)(s+1)=1+G(a,b,c).
 \label{eq:front-fused}
\end{equation}
This is the quadratic row $rs+r+s-a^2-b^2-c^2-c=0$; it requires no
multiplication gate.  Compose the previously supplied odd-square radix
restriction by deleting $q=B+o+1$ and its private guard, and inserting
\begin{equation}
 q=(2z+1)^2.
 \label{eq:front-odd}
\end{equation}
The two remaining guards---the strict squared Pell cap and the common
normalized remainder norm---are unchanged.  All other source rows,
all eighteen power requests, and all endpoint and Horn-flow relations
are retained.  The old coordinate $o$ occurs only in its deleted bound
row and can therefore be restored freely at the end.

The starting point is the direct component of
\path{ConZFC_Common_Remainders_Radix47_2026-10-25.zip}, at repository
snapshot \path{3ecd6a6cf53d47519228328567f45fb7ee9781d9}.
Its \path{NORMALIZED_REMAINDERS_PROOF.md}, together with the retained
\path{ORIENTATION_FIRST_PROOF.md} and
\path{STRICT_BUDGET_PROOF.md}, supplies the following interface
\emph{before} the removed front signs are used.  An integer $W\geq258$
and a signed Pell center $A$ satisfy
\begin{equation}
 |A|>W^W,\qquad |n|,|k|,|m|<W
 \label{eq:front-frame}
\end{equation}
for every requested base, index, and output.  The corresponding principal
Pell coordinates are canonical at the signed nominal index $k$.
The frame uses squares of the proposed atoms, and the squared semantic
anchor with bounded pooled errors orients both Pell coordinates.
Consequently those conclusions do not presuppose the signs of $F,B,q$
or of the independent CRT representatives.  We retain that inherited
orientation theorem and extend its subsequent power argument to signed
bases explicitly below.

\subsection{The signed-base extension}

\begin{lemma}[Signed bases]
\label{lem:front-signed}
Under the inherited interface \eqref{eq:front-frame}, a power request
with $2\leq |n|<W$ forces $k\geq0$ and $m=n^k$.  The assertion holds
for either sign of $n$ and either sign of $A$.
\end{lemma}
\begin{proof}
Set $T=2An-n^2-1$.  The actual output row is
\[
 m\equiv X_k(A)+(n-A)Y_k(A)\pmod T.
\]
The canonical signed pairs have
$X_{-j}(A)=X_j(A)$ and $Y_{-j}(A)=-Y_j(A)$.
Furthermore
\[
 |T|\geq2|A||n|-n^2-1
       \geq4|A|-(W-1)^2-1>W^W+W,
\]
where the last inequality follows from $W\geq258$ and $|A|>W^W$.
The absolute value on $A$ is essential: no positive-center premise has
been introduced.  Since $T\equiv-1\pmod n$, $n$ is invertible modulo $T$.

Substituting $n-A$ into the Pell identity gives
$(n-A)^2\equiv A^2-1\pmod T$.  The canonical Pell addition identities
therefore give $m\equiv n^k\pmod T$ when $k\geq0$.  In that case
\[
 |m-n^k|\leq |m|+|n|^k<W+W^W<|T|,
\]
so the congruence is an equality.  If $k=-j<0$, the signed pair is the
inverse of the index-$j$ Pell unit; its norm is one.  Thus
$mn^j\equiv1\pmod T$.  Now $1\leq j\leq W-1$ and $|m|\leq W-1$, so
\[
 |mn^j-1|\leq(W-1)^W+1<|T|.
\]
It follows that $mn^j=1$, impossible in $\Z$ because $|n|^j\geq2$.
Hence $k\geq0$ and the first argument applies.
\end{proof}

For the separate base-one case, a known nonnegative index gives
$m\equiv1\pmod{2A-2}$; the same size argument gives $m=1$.
This case is needed because $F=0$ has not yet been excluded.  The known
index $t\geq16$ in the request at base $F^2+1$ consequently gives
\begin{equation}
 V=(F^2+1)^t\geq1
 \label{eq:front-V}
\end{equation}
for every initially possible integer $F$.

\subsection{Both front branches have the needed size bound}

The untouched ceiling relations are
\begin{equation}
 r+s=B-1,\qquad BF=V+s.
 \label{eq:front-ceiling}
\end{equation}
The right side of \eqref{eq:front-fused} is positive.  Since its two
factors are integers, it leaves precisely the alternatives
\[
 \text{(I) }r,s\geq0,
 \qquad
 \text{(II) }r,s\leq-2.
\]
In case I, $B\geq1$ and $BF=V+s\geq1$ force $F\geq1$.
Since $V=(F^2+1)^t>F^2$, the inequality $V\leq BF$ then gives
$B>F$.  In particular $B\geq2$ and
\[
 V\leq BF<B^2.
\]

In case II put $B=-b$, so $b\geq3$.  If $F=0$, the base-one case of
\eqref{eq:front-V} gives $V=1$, and $BF=V+s$ would give $s=-1$.
If $F\geq1$, combining \eqref{eq:front-ceiling} gives
\[
 r=(1-F)B+V-1>0,
\]
because $(1-F)B\geq0$ and $V>1$.  Both possibilities contradict
case II.  Therefore $F=-f$ for some $f\geq1$.  The same two ceiling
identities now give
\begin{equation}
 bf+2\leq V\leq(f+1)b-1<(f+1)b.
 \label{eq:front-negative-bound}
\end{equation}
For $f\geq1$ and $t\geq16$, $(f^2+1)^t>(f+1)^2$.
Comparing this with \eqref{eq:front-negative-bound} shows $b>f+1$,
and hence $V<b^2$.

Thus, in both branches, defining $f=|F|$ and $b=|B|$ gives
\begin{equation}
 f\geq1,\qquad b\geq2,\qquad V=(f^2+1)^t<b^2.
 \label{eq:front-uniform}
\end{equation}
This deduction has used no factorial fixed point or binomial-digit
interpretation.  That separation is what makes the later sign recovery
noncircular.

\subsection{Recovering the radix before resolving the branch}

Equation \eqref{eq:front-odd} first gives $q\geq1$, so $q+1\geq2$.
Each of the nine actual CRT representatives $a_i$ is an exponent in a
retained request at base $q+1$.  Lemma~\ref{lem:front-signed} therefore
makes each $a_i\geq0$ and authenticates its output.  Set
\[
 H_0=1+\sum_{i=1}^{9}a_i\geq1.
\]
These are actual representatives; several delivered coordinates are
shifted by one and need not themselves be nonnegative.

Apply the lemma next to the retained request $N=B^x$.  Its base has
$|B|\geq2$ by \eqref{eq:front-uniform}, so it forces $x\geq0$ even
while $q=1$ is still an initial possibility.  The unchanged cofinal beta
choice is
\[
 y=V(V-1)>f+2.
\]
In unshifted semantic coordinates, the retained endpoints are
\begin{equation}
 x=F+1+(y+1)u_0,
 \qquad
 x=1+\bigl(1+(t+1)y\bigr)u_t.
 \label{eq:front-endpoints}
\end{equation}
The second equality excludes $x=0$ because its bracket is an integer
greater than one.  If $x=1$, the first would say
$F+(y+1)u_0=0$, impossible since $0<|F|<y+1$.  Consequently $x\geq2$,
and the second equality then forces $u_t\geq1$.  Thus
\begin{equation}
 x\geq2+(t+1)y\geq8.
 \label{eq:front-x-bound}
\end{equation}
The coordinate translations used in the emitted polynomial do not alter
these semantic identities; their explicit shifts are recorded separately.

The request $q=(BH_0)^x$ now has $|BH_0|\geq2$ and is authenticated
by the lemma.  Since its output is positive, $B<0$ would force $x$ even.
Therefore in either branch
\begin{equation}
 N=b^x>0,\qquad q=b^xH_0^x\geq b^x,\qquad N\mid q.
 \label{eq:front-radix-bound}
\end{equation}
In particular $q>1$.  All denominator requests are now genuine positive
powers of $q$.  Write
\[
 M=q^x,\qquad D_0=q^t,\qquad D_1=q^{t+1},\qquad D_y=q^{y+1}.
\]
By \eqref{eq:front-x-bound}, these three denominators divide $M$.
Nothing in this argument has assumed $F=t!$ or the sign of $B$.

\subsection{The exact factorial digit chooses the positive branch}

For $m\geq1$, $t\geq1$ and $0\leq j\leq t$, the elementary estimate
\begin{equation}
 \binom{m+t}{j}\leq(m+1)^t
 \label{eq:front-binomial-bound}
\end{equation}
holds.  One proof is induction on $t$: Pascal's identity bounds each
nonedge entry in the next row by $2(m+1)^t\leq(m+1)^{t+1}$; the new
edge is bounded by
$\prod_{i=1}^{t+1}(1+m/i)\leq(m+1)^{t+1}$.
Apply this with $m=F^2\geq1$.  Every coefficient through position $t$
of $(q+1)^{F^2+t}$ is at most $V<b^2$.  Its true low remainder at
position $t$ is therefore
\[
 r^0=\sum_{j=0}^{t-1}\binom{F^2+t}{j}q^j,
 \qquad
 0\leq r^0<\frac{V D_0}{q-1}<\frac{D_0}{32}.
\]
For the last inequality, \eqref{eq:front-radix-bound} and
\eqref{eq:front-x-bound} give
\[
 \frac{V}{q-1}<\frac{b^2}{q-1}
 <\frac{2}{b^{x-2}}\leq\frac1{32}.
\]
Thus the displayed sum really is the ordinary remainder modulo $D_0$;
the polynomial tail is divisible by $D_0$.

Here is the precise role of the inherited common norm.  Its normalized
digit rows have $D_i v_i=M$ and
$v_iU_i=M Q_i+R_i$.  Since $D_i,M>0$, $v_i$ is a positive integer.
The digit row makes $v_i\mid R_i$, so $r_i=R_i/v_i$ is an integer.
The norm equation
\begin{equation}
 M^2=1+4\sum_{i=1}^{11}R_i^2+G(\alpha,\beta,\gamma)
 \label{eq:front-norm}
\end{equation}
gives $|R_i|<M/2$, hence $|r_i|<D_i/2$.
For the factorial row, $r_i\equiv r^0\pmod{D_0}$.
Together with $0\leq r^0<D_0/32$, this forces $r_i=r^0$.
No sign of $B$ is used in this recovery.

The exact factorial digit row is now
\[
 (q+1)^{F^2+t}=D_0(B+qh)+r^0.
\]
Divide by $D_0$ after subtracting $r^0$ and reduce modulo $q$ to obtain
\begin{equation}
 B\equiv C\pmod q,
 \qquad C=\binom{F^2+t}{t}>0.
 \label{eq:front-final-digit}
\end{equation}
Unlike the old proof, this argument does not yet assume $0\leq B<q$.
Instead, \eqref{eq:front-uniform} and the coefficient estimate give
\[
 |B-C|\leq b+C\leq b+V<b+b^2<q,
\]
the last inequality following from $q\geq b^x$, $b\geq2$ and $x\geq8$.
Equation \eqref{eq:front-final-digit} therefore gives $B=C>0$.
Case II is impossible.  We have recovered $r,s\geq0$, $F\geq1$,
and $q>B$.  Restore the old private slack by $o=q-B-1\geq0$ and fill
all three deleted ternary guards using surjectivity of $G$.

At this point the previous Common Remainders source has an integer
solution.  Its exact factorial fixed-point theorem gives $F=t!$.
Its remaining ten low-digit estimates, prime-power CRT range decoding,
and Horn-flow interpretation therefore apply with their original
premises.  The new argument has restored those premises in the order
needed: size bound, radix, normalized remainder, factorial digit, front
signs, then factoriality.

\subsection{Completeness and the parity restriction}

For the converse, begin with a finite C222 refutation and use the
retained cofinal arithmetic construction.  The odd-square radix
restriction is inherited from the separate Centered Residues proof,
included as \texttt{INHERITED\_ODD\_SQUARE\_PROOF.md}; its compatibility
with the present construction can be seen explicitly.

After factoriality, $F=t!$ is even.  Each $i\leq t$ divides $F$ and
$F^2/i$ is even, so
\[
 B=\binom{F^2+t}{t}=\prod_{i=1}^t\left(1+\frac{F^2}{i}\right)
\]
is odd.  Also $y=V(V-1)$ is even, so every beta CRT modulus is odd.
The CRT progression for $x$ consequently has odd period, allowing $x$
to be chosen even and as large as required.  Now $N=B^x$ and $N^2$
are odd.  The independent amplified CRT representatives can be shifted
by multiples of $N^2$ without changing their required residues.
Adjusting a third representative's parity makes $H_0$ odd while leaving
the two representatives used for the dominant anchor's parity free to
retain opposite parities.  The dominant lift can still be made large
enough for all strict bounds.  Thus $BH_0$ is odd and $x$ is even;
$q=(BH_0)^x$ is the square of an odd integer and
\eqref{eq:front-odd} has an integer witness $z$.

The genuine old front has $r,s\geq0$, so $rs+r+s\geq0$ has a
$G$ representation and satisfies \eqref{eq:front-fused}.  The old
cofinal low-remainder estimates give $0\leq r_i^0<D_i/32$ for all
eleven digits.  Setting $R_i=(M/D_i)r_i^0$ gives $R_i<M/32$ and
\[
 4\sum_{i=1}^{11}R_i^2
 <\frac{44}{1024}M^2<M^2.
\]
Hence the nonnegative integer
$M^2-1-4\sum_iR_i^2$ has a $G$ representation, preserving
\eqref{eq:front-norm}.  Every other retained row is the old row with
its cofinal witness.  This constructs a zero of the new raw source.

The raw modification therefore preserves \emph{existence} of integer
zeros in both directions.  It is not a bijection on all arbitrary old
private-witness tuples.  The later monic projections and integral
translations, in contrast, have explicit reverse-restoration maps.

\subsection{Exact endpoint and the logical calibration chain}

\begin{theorem}
\label{thm:direct248}
Let $P_{248}\in\Z[x_0,\ldots,x_{247}]$ be the sum of the squares of
the 182 delivered residuals in the principal profile.  With the
retained direct-C222 arithmetic and logical calibration,
\[
 \mathrm{ZFC}\vdash
 \left(\exists x\in\Z^{248}\;P_{248}(x)=0\right)
 \longleftrightarrow
 \text{there exists a finite ZFC refutation}.
\]
Consequently ZFC proves that the absence of integer zeros of $P_{248}$
is equivalent to $\operatorname{Con}(\mathrm{ZFC})$.
The same conclusions hold for the two delivered 246- and 245-variable
quartics.
\end{theorem}
\begin{proof}
Over the integers a finite sum of squares vanishes exactly when every
residual vanishes.  The supplied integral translations are bijections;
each monic elimination has an integer polynomial restoration, replayed
in reverse order.  The independently checked Horn factoring identities
preserve the source's solution existence.  The preceding soundness and
completeness arguments identify zeros of the new source with zeros of
the Common Remainders source.  The retained arithmetic proof decodes
those zeros into finite C222 refutations, and constructs arithmetic
witnesses from such refutations.

The last two logical bridges are the identified repository theorems
\path{KEndpoint.c222_iff_zf_refutation} and
\path{ConstructibleChoice.zf_refutation_iff_zfc_refutation}.
Thus the complete chain is: explicit quartic zero, new raw arithmetic
witness, old Common Remainders witness, finite C222 refutation, finite
ZF refutation, and finite ZFC refutation.  Every arrow used here is
either the new arithmetic argument above, a finite exact algebraic
certificate, or the named inherited calibration.  The first-order
arithmetical statement that no finite syntactic ZFC refutation exists
is $\operatorname{Con}(\mathrm{ZFC})$; negating the existence
equivalence gives the stated result.
\end{proof}

The package retains all 88 files in the parent's arithmetic and proof
closure, pinned byte-for-byte.  Its new finite verifier imports the
retained independent algebra checker, not the producing eliminator.  It
reconstructs the complete raw mutation, checks 18 Horn-factor edits,
replays every monic restoration and translation, and expands each full
sum of squares.  The principal profile has 203 monic elimination steps
and 70 translated coordinates.  Independent runs under ordinary Python
and under \texttt{python -O} pass, including five rejected deliberate
mutations; a clean rebuild reproduces all 22 generated JSON files
byte-for-byte.  These are algebraic checks of the emitted artifacts.
The unbounded arithmetic proof remains an ordinary mathematical proof,
with a separate independent review; no new end-to-end Lean elaboration
or proof of consistency is being claimed.
